% problem-000149 8 重氮试剂合成应用
\section*{8. 重氮试剂合成应用}
\textit{下列为分问。}

\subsection*{8-1}
重氮甲烷 $\left( \mathrm { C H } _ { 2 } \mathrm { N } _ { 2 } \right)$ 可以作为甲基化试剂、卡宾前体、亲核试剂以及偶极⼦等，⼴泛应⽤在有机合成中。由于 $\mathrm { C H _ { 2 } N _ { 2 } }$ 存在毒性⾼和安全（沸点 $^ { - 2 3 ^ { \circ } \mathrm { C } }$ ，易燃易爆）等问题，科学家开发了其替代试剂：三甲基硅基重氮甲烷 $\mathrm { ( M e _ { 3 } S i C H N _ { 2 } }$ ，沸点 $9 6 ~ ^ { \circ } \mathrm { C } )$

\subsection*{8-1}
-1 画出 $\mathrm { M e _ { 3 } S i C H N _ { 2 } }$ 的两个稳定共振式。

\subsection*{8-1}
-2 $\mathrm { M e _ { 3 } S i C H N _ { 2 } }$ 常作为⼀碳合成单元应⽤在增碳反应中。写出如下反应产物的结构简式，并写出⽤结构简式表述的反应机理（提⽰： $\mathrm { S c ( O T f ) _ { 3 } }$ 是路易斯酸，三甲基硅基可以进⾏[1,3]-σ迁移）。

（图：）

\subsection*{8-1}
-3 在天然产物 amathaspiramide A 的全合成中利⽤了 $\mathrm { M e _ { 3 } S i C H N _ { 2 } }$ 参与的偶极环加成反应。画出下图所⽰全合成路线中化合物C、D和F的结构简式（不考虑产物的⽴体化学）。

（图：）

\subsection*{8-2}
烯烃复分解反应（烯烃交换反应）能够以两个烯烃为原料合成新的烯烃，⾰新了分⼦合成的设计逻辑。近年来化学家⼜发展了烯烃和酮的复分解反应。依据烯烃复分解反应的转换⽅式推测如下反应的主要产物G，并给出该转化过程中两个关键中间体的结构简式。

（图：）
